\documentclass[tikz,border=8pt]{standalone}
\usepackage{fontspec}
\usetikzlibrary{arrows.meta,decorations.markings}

\setmainfont{Arial}
\setsansfont{Arial}

\definecolor{pagebg}{HTML}{F8FAFC}
\definecolor{border}{HTML}{CBD5E1}
\definecolor{textmain}{HTML}{0F172A}
\definecolor{textmuted}{HTML}{475569}
\definecolor{north}{HTML}{2563EB}
\definecolor{northfill}{HTML}{DBEAFE}
\definecolor{south}{HTML}{DC2626}
\definecolor{southfill}{HTML}{FEE2E2}
\definecolor{field}{HTML}{0F766E}
\definecolor{fieldlight}{HTML}{14B8A6}

\begin{document}
\begin{tikzpicture}[
  x=1mm,y=1mm,line cap=round,line join=round,font=\sffamily,
  fieldarrow/.style={
    postaction={
      decorate,
      decoration={markings,mark=at position 0.52 with {\arrow{Latex[length=2.7mm,width=1.9mm]}}}
    }
  }
]
  \fill[pagebg] (0,0) rectangle (360,230);
  \draw[fill=white,draw=border,line width=0.45mm,rounded corners=4mm]
    (6,6) rectangle (354,224);

  \node[font=\sffamily\bfseries\fontsize{18}{22}\selectfont,text=textmain]
    at (180,207) {Линии магнитной индукции постоянных магнитов};

  \node[font=\sffamily\bfseries\fontsize{13}{16}\selectfont,text=textmain]
    at (95,184) {Полосовой магнит};
  \node[font=\sffamily\bfseries\fontsize{13}{16}\selectfont,text=textmain]
    at (265,184) {Подковообразный магнит};

  % Bar magnet.
  \draw[fill={rgb,255:red,21;green,169;blue,226},draw=north,line width=0.55mm]
    (48,105) rectangle (95,130);
  \draw[fill={rgb,255:red,255;green,82;blue,82},draw=south,line width=0.55mm]
    (95,105) rectangle (142,130);

  % Outside field lines for a bar magnet, matched to the reference:
  % lines leave N on the left, arc around the magnet, and enter S on the right.
  \draw[draw=field,line width=0.5mm,fieldarrow]
    (46,130) .. controls (20,178) and (170,178) .. (144,130);
  \draw[draw=field,line width=0.5mm,fieldarrow]
    (47,126) .. controls (28,160) and (162,160) .. (143,126);
  \draw[draw=field,line width=0.5mm,fieldarrow]
    (47,121) .. controls (36,146) and (154,146) .. (143,121);
  \draw[draw=field,line width=0.5mm,fieldarrow]
    (46,105) .. controls (20,57) and (170,57) .. (144,105);
  \draw[draw=field,line width=0.5mm,fieldarrow]
    (47,109) .. controls (28,75) and (162,75) .. (143,109);
  \draw[draw=field,line width=0.5mm,fieldarrow]
    (47,114) .. controls (36,89) and (154,89) .. (143,114);

  % Axial outside line: away from N on the left and into S on the right.
  \draw[-{Latex[length=2.7mm,width=1.9mm]},draw=field,line width=0.5mm]
    (48,117.5) -- (0,117.5);
  \draw[-{Latex[length=2.7mm,width=1.9mm]},draw=field,line width=0.5mm]
    (190,117.5) -- (142,117.5);

  % Inside the magnet the closed field lines return from S to N.
  \foreach \y in {112,117.5,123} {
    \draw[-{Latex[length=2.5mm,width=1.8mm]},draw=fieldlight,line width=0.48mm]
      (134,\y) -- (56,\y);
  }
  \node[font=\sffamily\bfseries\fontsize{13}{16}\selectfont,text=textmain,fill={rgb,255:red,21;green,169;blue,226},inner sep=0.8pt]
    at (60,117.5) {$N$};
  \node[font=\sffamily\bfseries\fontsize{13}{16}\selectfont,text=textmain,fill={rgb,255:red,255;green,82;blue,82},inner sep=0.8pt]
    at (130,117.5) {$S$};

  \node[font=\sffamily\fontsize{10}{12}\selectfont,text=textmuted,align=center]
    at (95,35) {Вне магнита линии идут\\от $N$ к $S$ и замыкаются внутри};

  % Horseshoe magnet: one U-shaped body with two facing poles.
  \draw[fill=northfill,draw=north,line width=0.65mm,rounded corners=2mm]
    (215,132) rectangle (305,154);
  \draw[fill=southfill,draw=south,line width=0.65mm,rounded corners=2mm]
    (215,72) rectangle (305,94);
  \draw[fill=pagebg,draw=border,line width=0.65mm,rounded corners=2mm]
    (215,72) rectangle (238,154);
  \draw[fill=white,draw=border,line width=0.45mm,rounded corners=1.5mm]
    (241,96) rectangle (303,130);

  \node[font=\sffamily\bfseries\fontsize{14}{17}\selectfont,text=north] at (288,143) {$N$};
  \node[font=\sffamily\bfseries\fontsize{14}{17}\selectfont,text=south] at (288,83) {$S$};

  % Nearly uniform field in the gap, from N to S.
  \foreach \x in {250,260,270,280,290} {
    \draw[-{Latex[length=3.2mm,width=2.2mm]},draw=field,line width=0.65mm]
      (\x,129) -- (\x,97);
  }

  % Outer closed field hints around the open side.
  \draw[-{Latex[length=3.2mm,width=2.3mm]},draw=field,line width=0.65mm]
    (307,145) .. controls (333,140) and (333,86) .. (307,82);
  \draw[-{Latex[length=3.2mm,width=2.3mm]},draw=field,line width=0.6mm]
    (306,153) .. controls (346,148) and (346,78) .. (306,73);
  \draw[-{Latex[length=3mm,width=2.1mm]},draw=fieldlight,line width=0.5mm]
    (226,83) -- (226,142);

  \node[font=\sffamily\fontsize{10}{12}\selectfont,text=textmuted,align=center]
    at (265,35) {В зазоре линии почти параллельны:\\поле сильнее и ближе к однородному};
\end{tikzpicture}
\end{document}
